\exercisesection{正则代数中的纯量扩张}

\begin{enumerate}
\def\labelenumi{\arabic{enumi})}
\tightlist
\item
  设 \(k\) 为特征 \(p > 2\) 的非完全域，\(a\) 为 \(k - k^p\) 的一个元素，且 \(A = k[\mathrm{X},\mathrm{Y}] / (\mathrm{X}^2 -\mathrm{Y}^p +a)\)。
\end{enumerate}

\begin{enumerate}
\def\labelenumi{\alph{enumi})}
\item
  证明环 A 是正则的（可以证明 A 在其分式域 \(K = k(Y)[X]/(X^{2} - Y^{p} + a)\) 中整闭）。
\item
  证明 k 在 K 中是代数闭的。
\item
  设 \(k'\) 为 \(k(a^{1/p})\)，它是 \(k\) 的一个扩张，证明环 \(A \otimes_k k'\) 不是正规的。
\end{enumerate}

\begin{enumerate}
\def\labelenumi{\arabic{enumi})}
\setcounter{enumi}{1}
\tightlist
\item
  设 \(k\) 为域，A 为一个本质有限型 \(k\)-代数，B 为一个 \(k\)-代数，\(p\) 为整数。
\end{enumerate}

\begin{enumerate}
\def\labelenumi{\alph{enumi})}
\item
  若 A 与 B 都满足性质 \(S_{p}\)（§ 1, 习题 8），则 \(A \otimes_{k} B\) 亦然（按命题 4 的证明方式论证，并利用该处 f)）。
\item
  设 K 为 k 的一个扩张；为使 \(\mathrm{A}_{(\mathbb{K})}\) 满足 \(S_{p}\)，必须且只需 A 满足之。
\end{enumerate}

\begin{enumerate}
\def\labelenumi{\arabic{enumi})}
\setcounter{enumi}{2}
\tightlist
\item
  设 A 为 Noether 环，\(\rho : A \to B\) 为环同态。若 B 是 Noether 环且是平坦 A-模，并且对每个 \(\mathfrak{p} \in \operatorname{Spec}(A)\)，\(\kappa(\mathfrak{p}) \otimes_{A} B\) 都是绝对正则的，则称 A-代数 B 是绝对正则的。
\end{enumerate}

\begin{enumerate}
\def\labelenumi{\alph{enumi})}
\item
  设 T 为 B 的一个乘性子集，S 为 A 的一个使 \(\rho(S) \subset T\) 的子集。若 B 是正则 A-代数，则 \(T^{1}B\) 是 \(S^{-1}A\) 上的正则代数。
\item
  为使 Noether A-代数 B 是绝对正则的，必须且只需对 B 的每个极大理想 n，\(A_{\rho^{-1}(n)}\)-代数 \(B_{n}\) 都是绝对正则的。
\item
  若 C 是绝对正则 B-代数且 B 是绝对正则 A-代数，则 C 是绝对正则 A-代数。
\item
  若 B 是 Noether 环且 C 是忠实平坦 B-模又是绝对正则 A-代数，则 B 是绝对正则 A-代数。
\item
  设 B 为绝对正则 A-代数，A′ 为一个本质有限型 A-代数。则 A′-代数 A′ ⊗A B 是绝对正则的。
\end{enumerate}

\begin{enumerate}
\def\labelenumi{\arabic{enumi})}
\setcounter{enumi}{3}
\item
  设 A 为 Noether 环，B 为一个忠实平坦的绝对正则 A-代数（习题 3）。为使 A 是正则（相应地，正规、Macaulay、Gorenstein）环，必须且只需 B 亦然。
\item
  若 A 是 Noether 环，并且对 A 的每个素理想 p，局部环 \(\Lambda_{p}\) 的完备化都是绝对正则的 \(A_{p}\)-代数（习题 3），则称 A 为 Grothendieck 环。
\end{enumerate}

\begin{enumerate}
\def\labelenumi{\alph{enumi})}
\item
  证明 Grothendieck 环的每个分式环与每个商环都是 Grothendieck 环。
\item
  设 A 为 Noether 环，B 为一个忠实平坦的绝对正则 A-代数。若 B 是 Grothendieck 环，证明 A 亦然（设 p 为 A 的素理想，q 为 B 的位于 p 之上的素理想，将利用 III, § 5, 第 4 节, 命题 4 证明 \(\widehat{B}_{q}\) 是 \(\widehat{\Lambda_{p}}\)-模且忠实平坦；然后应用习题 3）。
\end{enumerate}

\begin{enumerate}
\def\labelenumi{\arabic{enumi})}
\setcounter{enumi}{5}
\item
  设 A 为 Grothendieck 环，I 为 A 的理想，\(\widehat{\mathsf{A}}\) 为 A 关于 I-进拓扑的分离完备化。证明 \(\widehat{\mathsf{A}}\) 是绝对正则 A-代数。（设 \(\mathfrak{n}\) 为 \(\widehat{\mathsf{A}}\) 的一个极大理想，\(\mathfrak{m}\) 为其在 A 中的逆像；由习题 3 推出 \(\widehat{\Lambda}_{\mathfrak{n}}\) 是绝对正则的 \(\mathrm{A}_{\mathfrak{m}}\)-代数。）
\item
  \begin{enumerate}
  \def\labelenumii{\alph{enumii})}
  \tightlist
  \item
    证明特征 0 的整 Dedekind 环是 Grothendieck 环。
  \end{enumerate}
\end{enumerate}

\begin{enumerate}
\def\labelenumi{\alph{enumi})}
\setcounter{enumi}{1}
\tightlist
\item
  设 \(k\) 为特征 \(p > 0\) 的域，且 \(k\) 在 \(k^p\) 上是无限维的，\(r\) 为整数 \(\geqslant 1\)，A 为 \(k[[\mathrm{X}_1, \ldots, \mathrm{X}_r]]\) 中由系数生成 \(k^p\) 上有限维向量空间的形式幂级数所组成的子环。证明 A 是维数 \(r\) 的正则局部环，其完备化与 \(k[[\mathrm{X}_1, \ldots, \mathrm{X}_r]]\) 等同，但它不是 Grothendieck 环（利用 III, § 3 习题 17）。
\end{enumerate}

